% Topic 3: Simultaneous Linear Equations with Two Variables
% Part of the Matrix Fundamentals section

\subsection{Simultaneous Linear Equations with Two Variables}

Before using matrices to solve equations, it is important to review what a system of simultaneous linear equations actually represents. A linear equation in two variables, such as $ax + by = c$, geometrically represents a straight line on a Cartesian plane.

When we solve two simultaneous linear equations, we are looking for the $(x, y)$ coordinates that satisfy both equations at the same time. Geometrically, this means finding the point(s) where the two lines intersect.

\subsubsection*{Three Possible Outcomes}

For any system of two linear equations:
\[
\begin{cases}
a_1 x + b_1 y = c_1 \\
a_2 x + b_2 y = c_2
\end{cases}
\]
There are three possible geometric configurations and algebraic outcomes:

\begin{itemize}[leftmargin=*]
  \item \textbf{Unique Solution:} The lines have different gradients and intersect at exactly \textbf{one point}.
  \item \textbf{No Solution:} The lines have the same gradient but different $y$-intercepts. They are \textbf{parallel} and never intersect. Algebraically, solving them yields a contradiction (e.g., $0 = 5$).
  \item \textbf{Infinitely Many Solutions:} The lines have the same gradient and the same $y$-intercept. They are \textbf{coincident} (the exact same line). Algebraically, solving them yields a tautology (e.g., $0 = 0$).
\end{itemize}

\bigskip

% ── Worked Example: Solving Systems Algebraically ──
\begin{problem}[Algebraic and Geometric Interpretations]
For each of the following systems of equations, solve algebraically using either substitution or elimination, and describe the geometric relationship between the two lines.

\textbf{(a)} \\
$2x + 3y = 8$ \\
$5x - y = 3$

\medskip
\textbf{(b)} \\
$x + 2y = 4$ \\
$2x + 4y = 10$
\end{problem}

\begin{solution}
\textbf{(a)} \\
Let the equations be:
\begin{align*}
2x + 3y &= 8 \quad \text{--- (1)} \\
5x - y &= 3 \quad \text{--- (2)}
\end{align*}
Using substitution, rearrange equation (2) to make $y$ the subject: \\
$y = 5x - 3$

Substitute this into equation (1):
\begin{align*}
2x + 3(5x - 3) &= 8 \\
2x + 15x - 9 &= 8 \\
17x &= 17 \\
x &= 1
\end{align*}
Substitute $x = 1$ back into the rearranged equation (2): \\
$y = 5(1) - 3 = 2$

\textbf{Algebraic Solution:} $x = 1, y = 2$. \\
\textbf{Geometric Relationship:} The two lines have a unique solution. They intersect at the single point $(1, 2)$.

\medskip
\textbf{(b)} \\
Let the equations be:
\begin{align*}
x + 2y &= 4 \quad \text{--- (1)} \\
2x + 4y &= 10 \quad \text{--- (2)}
\end{align*}
Using elimination, multiply equation (1) by 2:
\[
2x + 4y = 8 \quad \text{--- (3)}
\]
Subtract equation (3) from equation (2):
\begin{align*}
(2x + 4y) - (2x + 4y) &= 10 - 8 \\
0 &= 2
\end{align*}
This is a mathematical contradiction ($0$ can never equal $2$).

\textbf{Algebraic Solution:} No solution. \\
\textbf{Geometric Relationship:} The two equations represent distinct parallel lines. Because they are parallel, they never intersect.
\end{solution}

\begin{takeaways}
\item \textbf{Algebra meets Geometry:} Solving simultaneous equations is identical to finding the intersection points of lines on a graph.
\item \textbf{Contradictions reveal Parallel Lines:} If your algebraic steps lead to an impossible statement like $0 = 2$, it is not a mistake. It is the mathematical way of stating the lines are parallel.
\end{takeaways}
